% Section 01 - Introduction
\section{Introduction}

\subsection{À propos de cette application}

Ce logiciel est une application de bureau développée spécifiquement pour la Faculté de Médecine et Sciences Pharmaceutiques (FMSP). Il permet de générer automatiquement :

\begin{itemize}[leftmargin=*]
    \item \textbf{Relevés de notes} détaillés pour les étudiants, avec calcul automatique des moyennes
    \item \textbf{Attestations de réussite} pour les étudiants ayant validé leurs semestres (moyenne $\geq$ 10/20)
    \item \textbf{Documents avec codes QR} pour faciliter la vérification et l'authenticité
\end{itemize}

L'application fonctionne entièrement \textbf{hors ligne} et vos données restent strictement confidentielles sur votre ordinateur.

\subsection{Fonctionnalités principales}

\subsubsection{Génération de documents}

\begin{description}[leftmargin=3cm, style=nextline]
    \item[Relevés de notes]
    \begin{itemize}
        \item Import de fichiers Excel avec données des étudiants
        \item Calcul automatique des moyennes par UE et moyenne générale
        \item Attribution automatique des mentions (Très Bien, Bien, Assez Bien, Passable, Ajourné)
        \item Gestion des crédits ECTS
        \item Support des sessions de rattrapage pour certains ECs
        \item QR code chiffré optionnel pour sécurité
        \item Génération par semestre ou multi-semestres
    \end{itemize}

    \item[Attestations de réussite]
    \begin{itemize}
        \item Validation automatique : seuls les étudiants avec moyenne $\geq$ 10/20 sont éligibles
        \item Statistiques en temps réel des étudiants admis/ajournés
        \item Thèmes entièrement personnalisables (couleurs, polices, mise en page)
        \item 6 préréglages de thèmes professionnels
        \item Support bilingue français/anglais
        \item QR code chiffré compact
    \end{itemize}

    \item[Codes QR sur documents]
    \begin{itemize}
        \item Ajout de QR codes sur PDFs existants (relevés, attestations, autres)
        \item QR code appliqué sur \textbf{TOUTES les pages} de chaque document
        \item Positionnement visuel précis (grille A4)
        \item Tailles personnalisables : petit, moyen, grand, ou taille personnalisée
        \item Niveaux de correction d'erreur ajustables (L, M, Q, H)
        \item Traitement en lot depuis Excel
    \end{itemize}
\end{description}

\subsubsection{Formats d'exportation}

L'application propose \textbf{3 formats d'exportation} pour tous les types de documents :

\begin{table}[h]
\centering
\begin{tabular}{|l|p{10cm}|}
\hline
\textbf{Format} & \textbf{Description} \\
\hline
\textbf{ZIP} & Archive compressée contenant tous les PDFs individuels. Format par défaut recommandé pour la distribution. \\
\hline
\textbf{Fichiers individuels} & Téléchargement séparé de chaque PDF (sans archive). Utile pour traiter un petit nombre de documents. \\
\hline
\textbf{PDF unique} & Tous les documents fusionnés en un seul fichier PDF. Idéal pour l'archivage centralisé. \\
\hline
\end{tabular}
\caption{Formats d'exportation disponibles}
\end{table}

\subsection{Les 9 menus de l'application}

L'interface de l'application est organisée en 9 menus accessibles depuis la barre latérale :

\begin{enumerate}[leftmargin=*]
    \item \menu{Générer les relevés} (\textit{Icône : FileSpreadsheet})
    \begin{itemize}
        \item Génération de relevés de notes à partir de fichiers Excel
        \item Composant : \texttt{ReleveGenerator}
    \end{itemize}

    \item \menu{Générer les attestations} (\textit{Icône : GraduationCap})
    \begin{itemize}
        \item Génération d'attestations de réussite
        \item Validation automatique des moyennes
        \item Composant : \texttt{AttestationGenerator}
    \end{itemize}

    \item \menu{Config des relevés} (\textit{Icône : Settings2})
    \begin{itemize}
        \item Configuration académique : semestres, UEs, ECs
        \item Gestion des coefficients et crédits
        \item Composant : \texttt{AcademicConfigManager}
    \end{itemize}

    \item \menu{Configurer les entêtes} (\textit{Icône : Cog})
    \begin{itemize}
        \item Paramètres de l'établissement (nom, logo, adresse)
        \item Configuration des informations affichées sur les documents
        \item Composant : \texttt{SettingsForm}
    \end{itemize}

    \item \menu{Historique des docs} (\textit{Icône : History})
    \begin{itemize}
        \item Consultation de l'historique des documents générés
        \item Recherche et filtrage
        \item Composant : \texttt{DocumentHistoryManager}
    \end{itemize}

    \item \menu{QR Codes sur PDF} (\textit{Icône : File})
    \begin{itemize}
        \item Ajout de QR codes sur PDFs existants
        \item Matching automatique avec données Excel
        \item Composant : \texttt{QrCodeOnPdf}
        \item \attention{Désactivé en mode démo}
    \end{itemize}

    \item \menu{Placement QR sur Documents} (\textit{Icône : QrCode})
    \begin{itemize}
        \item Placement avancé de QR codes avec mode manuel et Excel
        \item Personnalisation complète du QR code
        \item Composant : \texttt{QRCodeDocumentProcessor}
    \end{itemize}

    \item \menu{Aide} (\textit{Icône : HelpCircle})
    \begin{itemize}
        \item Documentation et support
        \item Composant : \texttt{HelpSupport}
    \end{itemize}

    \item \menu{Paramètres} (\textit{Icône : Settings})
    \begin{itemize}
        \item Configuration de l'application
        \item Gestion de la licence d'activation
        \item Composant : \texttt{SettingsPage}
    \end{itemize}
\end{enumerate}

\begin{figure}[h]
    \centering
    % \includegraphics[width=\textwidth]{images/interface_principale.png}
    \fbox{\parbox{0.9\textwidth}{\centering [CAPTURE D'ÉCRAN : INTERFACE PRINCIPALE AVEC SIDEBAR]\\ \vspace{0.3cm} Montrant les 9 menus dans la barre latérale gauche}}
    \caption{Interface principale de l'application}
\end{figure}

\subsection{Configuration requise}

\begin{tcolorbox}[colback=blue!5!white,colframe=blue!75!black,title=Configuration système minimale]
\begin{itemize}[leftmargin=*]
    \item \textbf{Système d'exploitation} : Windows 10/11, macOS 10.14+, ou Linux (Ubuntu 18.04+)
    \item \textbf{Processeur} : Intel Core i3 ou équivalent (2 GHz minimum)
    \item \textbf{Mémoire RAM} : 4 Go minimum (8 Go recommandé pour génération massive)
    \item \textbf{Espace disque} : 500 Mo pour l'application + espace pour documents générés
    \item \textbf{Résolution écran} : 1280x720 minimum (1920x1080 recommandé)
\end{itemize}
\end{tcolorbox}

\begin{tcolorbox}[colback=green!5!white,colframe=green!75!black,title=Configuration recommandée pour usage intensif]
\begin{itemize}[leftmargin=*]
    \item \textbf{Processeur} : Intel Core i5 ou supérieur (quad-core)
    \item \textbf{Mémoire RAM} : 8 Go ou plus
    \item \textbf{Disque dur} : SSD pour de meilleures performances
    \item \textbf{Générateur PDF} : Installation de Chrome ou Chromium (pour conversion HTML vers PDF)
\end{itemize}
\end{tcolorbox}

\subsection{Formats de fichiers supportés}

\subsubsection{Fichiers en entrée}

\begin{itemize}[leftmargin=*]
    \item \textbf{Données étudiants} : Excel (.xlsx, .xls) ou CSV
    \item \textbf{Documents PDF} : Tous fichiers PDF standards (pour ajout de QR codes)
    \item \textbf{Configurations} : Fichiers JSON (export/import de configurations)
    \item \textbf{Thèmes} : Fichiers JSON (export/import de thèmes)
\end{itemize}

\subsubsection{Fichiers générés en sortie}

\begin{itemize}[leftmargin=*]
    \item \textbf{Relevés de notes} : PDF (A4 portrait)
    \item \textbf{Attestations de réussite} : PDF (A4 portrait ou paysage selon thème)
    \item \textbf{Archives} : ZIP (compression optionnelle)
    \item \textbf{Historique} : Fichiers JSON de traçabilité
\end{itemize}

\subsection{Sécurité et confidentialité}

\begin{tcolorbox}[colback=yellow!10!white,colframe=orange!75!black,title=Protection des données]
\textbf{Garanties de confidentialité :}
\begin{itemize}[leftmargin=*]
    \item \textbf{Fonctionnement 100\% hors ligne} : Aucune connexion internet requise
    \item \textbf{Stockage local uniquement} : Toutes les données restent sur votre ordinateur
    \item \textbf{Aucun envoi de données} : Aucune information n'est transmise à l'extérieur
    \item \textbf{QR codes chiffrés} : Utilisation de AES-128-ECB avec clé basée sur matricule
    \item \textbf{Historique local} : Traçabilité stockée localement dans le navigateur
\end{itemize}

\textbf{Recommandations :}
\begin{itemize}[leftmargin=*]
    \item Effectuez des sauvegardes régulières de vos configurations
    \item Conservez les fichiers Excel sources en lieu sûr
    \item Exportez vos thèmes personnalisés pour éviter toute perte
    \item Sauvegardez l'historique des documents si nécessaire
\end{itemize}
\end{tcolorbox}

\subsection{Mode démo et licence}

L'application propose deux modes de fonctionnement :

\begin{description}[leftmargin=3cm]
    \item[Mode démo]
    \begin{itemize}
        \item Accès à la plupart des fonctionnalités
        \item Limitation : le menu \menu{QR Codes sur PDF} est désactivé
        \item Idéal pour tester l'application avant activation
    \end{itemize}

    \item[Mode activé]
    \begin{itemize}
        \item Accès complet à toutes les fonctionnalités
        \item Nécessite une clé de licence valide
        \item Configuration via le menu \menu{Paramètres}
    \end{itemize}
\end{description}

\newpage
